% M.Med 模板样式配置（尽量集中到一个文件，便于维护与对齐 Word 模板）

%------------------------ 基础宏包 ------------------------%
\RequirePackage{geometry}
\RequirePackage{setspace}
\RequirePackage{fontspec}
\RequirePackage{xeCJK}
\RequirePackage{xcolor}
\RequirePackage{colortbl} % 提供 \arrayrulecolor（封面横线颜色微调）
\RequirePackage{calc} % 支持带小数的长度计算（pandoc longtable 使用 \real{...}）
\RequirePackage{graphicx}
\RequirePackage{newunicodechar} % 将 Word 导出的少量 Unicode 符号映射到可编译形式
\RequirePackage{float}
\RequirePackage[section]{placeins} % 避免浮动体跨越章节
\RequirePackage{booktabs}
\RequirePackage{longtable}
\RequirePackage{tabularx}
\RequirePackage{multirow}
\RequirePackage{caption}
\RequirePackage{amsmath,amssymb,bm}
\RequirePackage{titlesec} % 精细控制标题：用于实现“编号/标题分行”的节标题
\RequirePackage{enumitem}
\RequirePackage{hyperref}
\RequirePackage{fancyhdr}
\RequirePackage{etoolbox}
\RequirePackage[
  backend=biber,
  style=gb7714-2015,
  sorting=none
]{biblatex}

%------------------------ 颜色（对齐 Word 的 MsBlue） ------------------------%
% 说明：
% - Word 常用蓝色：RGB(0,112,192)
% - 在 XeLaTeX + xcolor 下，RGB(0,112,192) 在 PDF 中可能被量化为 (0,111,192)；
%   为了让最终 PDF 提取/显示更贴近 Word 的观感，这里将 G 分量 +1。
\definecolor{MsBlue}{RGB}{0,113,192}
% 兼容用户口头习惯写法（MSBlue / MsBlue）
\colorlet{MSBlue}{MsBlue}

% pandoc 生成的 longtable 可能会设置 \LTcaptype{none}；这里提供一个空计数器避免报错
\newcounter{none}

% 常见符号：避免字体缺字导致编译警告/断字
\newunicodechar{⁹}{\textsuperscript{9}}
\newunicodechar{△}{\(\Delta\)}
\newunicodechar{☑}{\(\checkmark\)}
\newunicodechar{‐}{-}

%------------------------ 页面尺寸（A4） ------------------------%
% 说明：
% - 左/右：约 3.2cm（与 Word 模板一致）
% - 上/下：此处先采用 2.7cm（当前更贴近现有 docx->PDF 的版心与页数）；如需更高保真，可再迭代微调。
\geometry{left=3.2cm,right=3.2cm,top=2.7cm,bottom=2.7cm,headheight=14pt,headsep=10pt}

%------------------------ 字体 ------------------------%
% 英文：优先 Times New Roman；若不存在则回退到 TeX Gyre Termes
\IfFontExistsTF{Times New Roman}{
  \setmainfont{Times New Roman}
}{
  \setmainfont{TeX Gyre Termes}
}

% 中文：优先 宋体（macOS 的 Songti SC）；若不存在则回退到 STSong
\IfFontExistsTF{Songti SC}{
  \setCJKmainfont{Songti SC}
}{
  \setCJKmainfont{STSong}
}

\IfFontExistsTF{Heiti SC}{
  \setCJKsansfont{Heiti SC}
}{}

% 避免 xeCJK “Unknown CJK family \\CJKttdefault” 警告
\IfFontExistsTF{Songti SC}{
  \setCJKmonofont{Songti SC}
}{}

%------------------------ 段落与行距 ------------------------%
\setstretch{1.5}
\setlength{\parindent}{2em}
\setlength{\parskip}{0pt}
\RequirePackage{indentfirst}

%------------------------ 目录标题 ------------------------%
% Word 模板常见为“目 录”（中间留空）；这里保留空隙以贴近观感
\renewcommand{\contentsname}{目\hspace{1em}录}

%------------------------ 超链接 ------------------------%
% 说明：
% - 正文中的 cite 引用（citecolor）与 URL 链接（urlcolor）使用 MsBlue
% - 目录中的页码/标题链接（linkcolor）默认也为 MsBlue，但在目录输出时局部覆盖为黑色
%   → 在 main.tex 的 \tableofcontents 前后各加一行 \hypersetup 即可（见下方工具宏）
\hypersetup{
  colorlinks=true,
  hypertexnames=false, % 避免多次 \pagenumbering 导致 PDF 书签重复锚点（Object @page.* already defined）
  linkcolor=MsBlue,
  citecolor=MsBlue,
  urlcolor=MsBlue,
  CJKbookmarks=true
}

% 目录局部颜色控制宏（在 main.tex 中包裹 \tableofcontents 使用）
% 用法：
%   \mmedTOCBlackLinks
%   \tableofcontents
%   \mmedTOCRestoreLinks
\newcommand{\mmedTOCBlackLinks}{%
  \hypersetup{linkcolor=black}%
}
\newcommand{\mmedTOCRestoreLinks}{%
  \hypersetup{linkcolor=MsBlue}%
}

%------------------------ 图表编号（避免手写序号） ------------------------%
% 说明：
% - 统一采用“章号-序号”的形式（如 3-1），便于与常见医学论文/毕业论文格式保持一致
% - 显式使用 arabic，避免 ctex 章节标题使用中文数字时影响图表计数显示
\renewcommand{\thefigure}{\arabic{chapter}-\arabic{figure}}
\renewcommand{\thetable}{\arabic{chapter}-\arabic{table}}

%------------------------ 图表标题 ------------------------%
\captionsetup{
  font={small,stretch=1.0},
  labelfont=bf,
  labelsep=space,
  justification=centering,
  singlelinecheck=false
}

%------------------------ 列表（贴近医学论文常用风格） ------------------------%
\setlist[enumerate]{
  leftmargin=2.5em,
  itemsep=0.3em,
  topsep=0.3em
}
\setlist[itemize]{
  leftmargin=2.5em,
  itemsep=0.3em,
  topsep=0.3em
}

%------------------------ 页眉页脚 ------------------------%
\newcommand{\mmedHeaderFont}{\zihao{-5}}
\newcommand{\mmedHeaderTextMain}{硕士学位论文}

%------------------------ 摘要首页微调 ------------------------%
% 说明：
% - Word 原稿的“摘要首页”在页眉横线（headrule）下方通常会留出更大的空白，
%   以便与题名区块形成明确的视觉分隔。
% - 这里将该留白参数化，避免通过全局 geometry/headsep 影响所有页面。
\newlength{\mmedAbstractTopGap}
\setlength{\mmedAbstractTopGap}{0.08cm}

\fancypagestyle{mmed-main}{
  \fancyhf{}
  \fancyhead[C]{\mmedHeaderFont \mmedHeaderTextMain}
  \fancyfoot[C]{\thepage}
  \renewcommand{\headrulewidth}{0.4pt}
  \renewcommand{\footrulewidth}{0pt}
}

\fancypagestyle{mmed-abstract-zh}{
  \fancyhf{}
  \fancyhead[C]{\mmedHeaderFont 摘\hspace{0.5em}要}
  \fancyfoot[C]{\thepage}
  \renewcommand{\headrulewidth}{0.4pt}
  \renewcommand{\footrulewidth}{0pt}
}

\fancypagestyle{mmed-abstract-en}{
  \fancyhf{}
  \fancyhead[C]{\mmedHeaderFont ABSTRACT}
  \fancyfoot[C]{\thepage}
  \renewcommand{\headrulewidth}{0.4pt}
  \renewcommand{\footrulewidth}{0pt}
}

%------------------------ 章节标题（尽量贴近 Word） ------------------------%
\ctexset{
  chapter = {
    name = {第,章},
    number = \chinese{chapter},
    format = \centering\bfseries\zihao{3},
    % Word 常见：章号与章名分两行居中显示
    aftername = \par\vspace*{0.3em},
    beforeskip = 14pt,
    afterskip = 14pt
  },
  section = {
    format = \bfseries\zihao{4},
    beforeskip = 8pt,
    afterskip = 6pt
  },
  subsection = {
    format = \bfseries\zihao{4},
    beforeskip = 10pt,
    afterskip = 4pt
  }
}

%------------------------ 小节标题样式 ------------------------%
% 说明：
% - 不能通过 ctexset 的 number/aftername 直接插入 \par（会污染 mark，导致 fancyhdr 报错）
% - 用 titlesec 仅改变”渲染外观”，不影响目录/书签/页眉的标记逻辑
% - 序号与标题在同一行，序号后空一格
\titleformat{\section}[block]
  {\bfseries\zihao{4}}
  {\thesection}
  {1em}
  {}
\titlespacing*{\section}{0pt}{8pt}{6pt}

%------------------------ 工具宏 ------------------------%
\newcommand{\mmedTOCAddChapter}[1]{%
  \phantomsection%
  \addcontentsline{toc}{chapter}{#1}%
}

% 压缩“中文-西文/数字”之间的自动间距（主要用于封面填写字段，贴近 Word 原稿）。
% 说明：仅在 XeLaTeX + xeCJK 环境生效；其他引擎下安全降级为原样输出。
\newcommand{\mmedTightCJK}[1]{%
  \begingroup
    \ifcsname xeCJKsetup\endcsname
      \xeCJKsetup{CJKecglue=\hspace{0pt}}%
    \fi
    #1%
  \endgroup
}

% 带下划线的“可填写字段”（封面用）：空值则保持空白下划线
\newcommand{\mmedUnderlinedField}[2]{%
  \underline{\makebox[#2][c]{\ifstrempty{#1}{}{#1}}}%
}

% 封面用：带下划线的“填写字段”（左对齐，并在内容后自动补齐下划线）
% - 典型用法：\mmedUnderlinedFieldL{张三}{\linewidth}
\newcommand{\mmedUnderlinedFieldL}[2]{%
  \underline{\hbox to #2{\ifstrempty{#1}{}{#1}\hfil}}%
}

% 封面用（带左内边距）：用于底部“导师/专业/培养类型/提交日期”等字段，
% 避免在紧凑列间距下出现“内容贴边/视觉上像与标签重叠”的问题。
\newcommand{\mmedUnderlinedFieldLPad}[2]{%
  \underline{\hbox to #2{%
    \ifstrempty{#1}{}{%
      % 字段内容略向右缩进，避免“内容贴住标签”造成视觉拥挤
      % 继续增大间隙（在上一版基础上约 +80%），让底部字段更接近原始模板留白
      \hspace*{1.44em}\strut#1%
    }%
    \hfil%
  }}%
}

% 封面标题多行分隔符（默认仅换行；cover.tex 内会局部重定义为“换行+横线”以贴近 Word）
\newcommand{\mmedCoverLinebreak}{\\}

% 摘要区块：左侧粗体标签 + 正文
\newcommand{\mmedAbsLabel}[1]{\noindent\textbf{#1}\ }

% 中英双语图表说明（用于 docx 转换后的“表3-1 ... / Table3-1 ...”这类文本说明）
% 说明：
% - 旧版实现为“纯文本说明行”（不参与 table 计数），容易导致图表序号需要手工维护
% - 新版改为 captionof(table)，由 LaTeX 自动编号；可选传入 label 以支持超链接引用
% - 用法示例：
%   \mmedBiCaption[tab:demo]{纳入患者的一般临床资料（n=174）}{General clinical data of included patients (n=174)}
\newcommand{\mmedZhCaption}[2][]{%
  \begingroup
    \captionsetup{type=table}%
    \captionof{table}{#2}%
    \ifstrempty{#1}{}{\label{#1}}%
  \endgroup
}
\newcommand{\mmedBiCaption}[3][]{%
  \begingroup
    \captionsetup{type=table}%
    \captionof{table}{#2}%
    \ifstrempty{#1}{}{\label{#1}}%
    \caption*{Table~\thetable~#3}%
  \endgroup
}

%------------------------ 参考文献样式（可调参数） ------------------------%
% 说明：
% - 以下参数均有默认值，可在 main.tex 的 \printbibliography 前按需覆盖。
% - 字号：\renewcommand{\mmedBibFontCmd}{\zihao{-5}}
% - 行距：\renewcommand{\mmedBibStretch}{1.5}
% - 条目间距：\setlength{\mmedBibItemSep}{4pt}
% - 条目内段间距：\setlength{\mmedBibParsep}{0pt}

% 字号（默认小六 = 6.5pt，约为正文小五的 75%）
\newcommand{\mmedBibFontCmd}{\zihao{-6}}
% 行距倍数（默认 1.0 = 单倍行距）
\newcommand{\mmedBibStretch}{1.0}
% 条目间额外间距（biblatex \bibitemsep）
\newlength{\mmedBibItemSep}
\setlength{\mmedBibItemSep}{1pt}
% 条目内段落间距（biblatex \bibparsep）
\newlength{\mmedBibParsep}
\setlength{\mmedBibParsep}{0pt}

% 在参考文献列表开始时自动应用上述参数
\AtBeginBibliography{%
  \mmedBibFontCmd
  \setstretch{\mmedBibStretch}%
  \setlength\bibitemsep{\mmedBibItemSep}%
  \setlength\bibparsep{\mmedBibParsep}%
}

% 有 DOI 时不再重复显示 URL（避免"已有 DOI 又展示 https 链接"的冗余）
\AtEveryBibitem{%
  \iffieldundef{doi}{}{\clearfield{url}\clearfield{urldate}}%
}

%------------------------ 参考文献引用（位置与间距） ------------------------%
% 说明：
% - 推荐写法：把 \cite 放在中文标点之前，例如：……发病与死亡负担长期居于前列\cite{xxx}。
% - XeLaTeX + xeCJK 默认会在 CJK 与 ASCII（如 \cite 输出的 [1]）之间插入间距，导致“前列 [1]”；
%   这里在常用 cite 命令开头加入 \unskip，移除紧邻其前的空格/胶水，得到“前列[1]”。
\ifcsdef{cite}{\pretocmd{\cite}{\unskip}{}{}}{}
\ifcsdef{parencite}{\pretocmd{\parencite}{\unskip}{}{}}{}
\ifcsdef{textcite}{\pretocmd{\textcite}{\unskip}{}{}}{}
\ifcsdef{autocite}{\pretocmd{\autocite}{\unskip}{}{}}{}
\ifcsdef{footcite}{\pretocmd{\footcite}{\unskip}{}{}}{}
\ifcsdef{supercite}{\pretocmd{\supercite}{\unskip}{}{}}{}

%------------------------ 常用引用宏（正文写作更省心） ------------------------%
% 说明：用 \hyperref 让“图/表/节”等完整短语都可点击（而不仅仅是编号）
\newcommand{\figref}[1]{\hyperref[#1]{图~\ref*{#1}}}
\newcommand{\tabref}[1]{\hyperref[#1]{表~\ref*{#1}}}
\newcommand{\secref}[1]{\hyperref[#1]{\ref*{#1}节}}
\newcommand{\chapref}[1]{\hyperref[#1]{第~\ref*{#1}~章}}
